\begin{thebibliography}{99}

\bibitem{BellmanAstrom1970}
Bellman, R. and \AA{}str\"om, K.J. (1970).
On structural identifiability.
\emph{Mathematical Biosciences}, 7, 329--339.

\bibitem{Rothenberg1971}
Rothenberg, T.J. (1971).
Identification in parametric models.
\emph{Econometrica}, 39, 577--591.

\bibitem{Manski2003}
Manski, C.F. (2003).
\emph{Partial Identification of Probability Distributions}.
Springer.
\url{https://doi.org/10.1007/b97478}

\bibitem{GraceEtAl2025}
Grace, J.B. et al. (2025).
Causal effects versus causal mechanisms: two traditions with different requirements and contributions towards causal understanding.
\emph{Ecology Letters}, 28, e70029.

\bibitem{CorreiaEtAl2025}
Correia, H.E., Dee, L.E., and Ferraro, P.J. (2025).
Designing causal mediation analyses to quantify intermediary processes in ecology.
\emph{Biological Reviews}, 100, 1512--1533.

\bibitem{SiegelDee2025}
Siegel, K. and Dee, L.E. (2025).
Foundations and future directions for causal inference in ecological research.
\emph{Ecology Letters}, 28, e70053.

\bibitem{MacKenzieEtAl2002}
MacKenzie, D.I., Nichols, J.D., Lachman, G.B., Droege, S., Royle, J.A., and Langtimm, C.A. (2002).
Estimating site occupancy rates when detection probabilities are less than one.
\emph{Ecology}, 83, 2248--2255.
\url{https://doi.org/10.1890/0012-9658(2002)083[2248:ESORWD]2.0.CO;2}

\bibitem{RoyleLink2006}
Royle, J.A. and Link, W.A. (2006).
Generalized site occupancy models allowing for false positive and false negative errors.
\emph{Ecology}, 87, 835--841.
\url{https://doi.org/10.1890/0012-9658(2006)87[835:GSOMAF]2.0.CO;2}

\bibitem{Dietze2017}
Dietze, M.C. (2017).
Prediction in ecology: a first-principles framework.
\emph{Ecological Applications}, 27, 2048--2060.
\url{https://doi.org/10.1002/eap.1589}

\bibitem{DietzeEtAl2018}
Dietze, M.C. et al. (2018).
Iterative near-term ecological forecasting: Needs, opportunities, and challenges.
\emph{Proceedings of the National Academy of Sciences USA}, 115, 1424--1432.
\url{https://doi.org/10.1073/pnas.1710231115}

\bibitem{NicholsWilliams2006}
Nichols, J.D. and Williams, B.K. (2006).
Monitoring for conservation.
\emph{Trends in Ecology \& Evolution}, 21, 668--673.
\url{https://doi.org/10.1016/j.tree.2006.08.007}

\bibitem{LindenmayerLikens2009}
Lindenmayer, D.B. and Likens, G.E. (2009).
Adaptive monitoring: a new paradigm for long-term research and monitoring.
\emph{Trends in Ecology \& Evolution}, 24, 482--486.
\url{https://doi.org/10.1016/j.tree.2009.03.005}

\bibitem{GivanDeanGreig2003}
Givan, R., Dean, T., and Greig, M. (2003).
Equivalence notions and model minimization in Markov decision processes.
\emph{Artificial Intelligence}, 147, 163--223.
\url{https://doi.org/10.1016/S0004-3702(02)00376-4}

\bibitem{ChalonerVerdinelli1995}
Chaloner, K. and Verdinelli, I. (1995).
Bayesian experimental design: A review.
\emph{Statistical Science}, 10, 273--304.
\url{https://doi.org/10.1214/ss/1177009939}

\bibitem{VanlierEtAl2012}
Vanlier, J., Tiemann, C.A., Hilbers, P.A.J., and van Riel, N.A.W. (2012).
A Bayesian approach to targeted experiment design.
\emph{Bioinformatics}, 28, 1136--1142.
\url{https://doi.org/10.1093/bioinformatics/bts092}

\bibitem{AttiaEtAl2018}
Attia, A., Alexanderian, A., and Saibaba, A.K. (2018).
Goal-oriented optimal design of experiments for large-scale Bayesian linear inverse problems.
\emph{Inverse Problems}, 34, 095009.
\url{https://doi.org/10.1088/1361-6420/aad210}

\bibitem{ZhongEtAl2026}
Zhong, S., Shen, W., Catanach, T., and Huan, X. (2026).
Goal-Oriented Bayesian Optimal Experimental Design for Nonlinear Models Using Markov Chain Monte Carlo.
\emph{SIAM/ASA Journal on Uncertainty Quantification}, 14, 19--47.
\url{https://doi.org/10.1137/24M1649344}

\bibitem{CanessaEtAl2015}
Canessa, S., Guillera-Arroita, G., Lahoz-Monfort, J.J., Southwell, D.M., Armstrong, D.P., Chad\`es, I., Lacy, R.C., and Converse, S.J. (2015).
When do we need more data? A primer on calculating the value of information for applied ecologists.
\emph{Methods in Ecology and Evolution}, 6, 1219--1228.
\url{https://doi.org/10.1111/2041-210X.12423}

\bibitem{BeaumontEtAl2002}
Beaumont, M.A., Zhang, W., and Balding, D.J. (2002).
Approximate Bayesian computation in population genetics.
\emph{Genetics}, 162, 2025--2035.

\bibitem{RobertEtAl2011}
Robert, C.P., Cornuet, J.-M., Marin, J.-M., and Pillai, N.S. (2011).
Lack of confidence in approximate Bayesian computation model choice.
\emph{Proceedings of the National Academy of Sciences USA}, 108, 15112--15117.

\bibitem{BetiniEtAl2017}
Betini, G.S., Avgar, T., and Fryxell, J.M. (2017).
Why are we not evaluating multiple competing hypotheses in ecology and evolution?
\emph{Royal Society Open Science}, 4, 160756.

\bibitem{YancoEtAl2020}
Yanco, S.W., McDevitt, A., Trueman, C.N., Hartley, L., and Wunder, M.B. (2020).
A modern method of multiple working hypotheses to improve inference in ecology.
\emph{Royal Society Open Science}, 7, 200231.

\end{thebibliography}
